
\section*{6. Passage to the quotient when a quasi-section exists}
\phantomsection\label{sec:sga3-V-6}
\addcontentsline{toc}{section}{6. Passage to the quotient when a quasi-section exists}

We shall now prove a technical lemma that will be used in the proof of
the two theorems we have in view. Let \(S\) be a scheme and let
\SGAThreeDiagram[0.54\textwidth]
  {figures/exp5/p270_section6_groupoid_X.png}
  {The groupoid \(X_\ast\).}
  {fig:sga3-V-6-X-groupoid}
  {LNM 151, pp.~271--272}
be a \((\mathbf{Sch}/S)\)-groupoid. We call a \emph{quasi-section} of
the groupoid \(X_\ast\) any subscheme \(U\) of \(X_0\) for which the
following two conditions hold:
\begin{enumerate}
\item[(1)]\phantomsection\label{def:sga3-V-6-quasisection-1}%
The restriction \(v\) of \(d_1\) to \(d_0^{-1}(U)\) is a
  finite, locally free, surjective morphism from \(d_0^{-1}(U)\) onto
  \(X_0\).
\item[(2)]\phantomsection\label{def:sga3-V-6-quasisection-2}%
Every subset \(E\) of \(U\) consisting of points that are
  pairwise equivalent for the equivalence relation defined by
  \(X_\ast\) (\hyperref[par:sga3-V-3-e]{\S3.e)}) is contained in an affine open
  subset of \(U\).
\end{enumerate}

If \(U\) is a quasi-section of \(X_\ast\), the
\((\mathbf{Sch}/S)\)-groupoid
\SGAThreeDiagram[0.48\textwidth]
  {figures/exp5/p270_section6_groupoid_U.png}
  {The groupoid \(U_\ast\) induced by \(X_\ast\).}
  {fig:sga3-V-6-U-groupoid}
  {LNM 151, pp.~271--272}
induced by \(X_\ast\) and the inclusion \(U\hookrightarrow X_0\)
(see \hyperref[par:sga3-V-3-a]{\S3.a)}) satisfies the hypotheses of
Theorem~\hyperref[thm:sga3-V-4.1]{4.1}. Indeed, put
\[
  V=d_0^{-1}(U),
\]
and let \(u\) and \(v\) be the morphisms with source \(V\) induced by
\(d_0\) and \(d_1\), respectively:
\SGAThreeDiagram[0.43\textwidth]
  {figures/exp5/p270_section6_span_X0_V_U.png}
  {The morphisms \(v:V\to X_0\) and \(u:V\to U\).}
  {fig:sga3-V-6-span-X0-V-U}
  {LNM 151, pp.~271--272}

By \hyperref[par:sga3-V-3-b]{\S3.b)}, there is a Cartesian square
\SGAThreeDiagram[0.27\textwidth]
  {figures/exp5/p270_section6_cartesian_square.png}
  {The Cartesian square defining \(U_1\).}
  {fig:sga3-V-6-U1-cartesian}
  {LNM 151, pp.~271--272}
so \(u_1\) is surjective and finite locally free by \hyperref[def:sga3-V-6-quasisection-1]{condition~(1)}.
Together with \hyperref[def:sga3-V-6-quasisection-2]{condition~(2)}, \hyperref[def:sga3-V-6-quasisection-1]{condition~(1)} therefore ensures that the
groupoid \(U_\ast\) satisfies the hypotheses of
Theorem~\hyperref[thm:sga3-V-4.1]{4.1}. In particular,
\(\operatorname{Coker}(u_0,u_1)\) exists in \((\mathbf{Sch}/S)\).
Moreover, \(d_0\) has a section, so \(u\) is a universal effective
epimorphism (see Expos\'e IV, 1.12).\footnote{The French re-edition
mistakenly cites Expos\'e III, which has no \S1.12.}
Proposition~\hyperref[prop:sga3-V-3.1]{3.1}
then shows that \(\operatorname{Coker}(u_0,u_1)\) coincides with the
cokernel \(\operatorname{Coker}(v_0,v_1)\) of the groupoid \(V_\ast\):
\SGAThreeDiagram[0.48\textwidth]
  {figures/exp5/p270_section6_groupoid_V.png}
  {The groupoid \(V_\ast\).}
  {fig:sga3-V-6-V-groupoid}
  {LNM 151, pp.~271--272}
